\section{Implementations}\label{sec:implementations}

In this section we will be using the numerical methods discussed in Section~\ref{sec:numerical_methods} to implement physical models as discussed in Section~\ref{sec:physics}.

We start with a `0D' model of absorption and desorption. This model can be seen as a perfectly mixed system, where the concentration of the H2 gas is uniform over the system. The first adsorption/desorption model will be built on a density measure, where the density is seen as the percentage of hydrogen gas in the air. We will built upon this in the second model by using a concentration measure, where the concentration is seen as the amount of hydrogen gas in a certain volume. This second model will be the one dimensional application of diffusion, absorption, and convection using the finite difference method. This model assumes a pre-determined speed of flow for the gas through the system. The third model will be a two dimensional application of diffusion, absorption, and convection using the finite element method.

The first three models give a good view on the applications of the numerical methods. The underlying physics can be seen as a system where a battery is filled with normal gas, and the entry point is opened to let the hydrogen gas in. In the next models we will be looking to implement a more realistic model, where the flow of the gas is not pre-determined, but is instead calculated using the Navier-Stokes equations. Also, we want to say something about pressure, which we do by letting the gas flow through a porous medium completely filled with hydrogen gas. The pressure can then be approximated with the ideal gas law. Pressure plays a role in the absorption, so our model 4 is a pressure based absorption/desorption `0D' model. Model 5 is a one dimensional application of diffusion, absorption, and convection using the finite difference method, where the flow of the gas is calculated using the Navier-Stokes equations. Model 6 is a two dimensional application of diffusion, absorption, and convection using the finite element method, where the flow of the gas is calculated using the Navier-Stokes equations.